% GENERATED FILE. Edit canonical lessons, then run npm run generate:books.
% canonical-source-hash: fnv1a64:b1fe22eb1ce7790a
% canonical-lessons: ZH-C114-maidan, ZH-C114-fu, ZH-C114-xianjin, ZH-C114-shuaka, ZH-C114-zhaoqian

\chapter{Paying}
\label{ch:zh-paying}
\hlchaptermodality{114}

\begin{chapteropening}
\textbf{By the end of this chapter:} \emph{I can ask for the bill, pay in cash or by card, and get my change.}

Complete the last lesson of chapter 114: I can ask for the bill, pay in cash or by card, and get my change.
\end{chapteropening}

\section[mǎidān]{mǎidān — to pay the bill}

\label{lesson:ZH-C114-maidan}



\begin{quote}
Before the new one: say the Mandarin for to roast, then the Mandarin for to fry.
\end{quote}

\subsection*{You'll want to know: mǎidān}
\emph{mǎidān} — \textquotedblleft{}to pay the bill\textquotedblright{}.

\begin{quote}
\emph{fúwùyuán mǎidān} — Excuse me, the bill, please!
\end{quote}

Its characters come later; for now it is a word for the ear.

\begin{cousinweb}
\emph{mǎidān} is built from \zh{买} 'to buy' and dān, both inherited from Middle Chinese, \emph{meaX} and \emph{tan} (Guangyun, in Baxter's transcription); the dictionary glosses the second 'single'. Paying the bill is a modern sense, and how the pair came to it is not settled.
\end{cousinweb}

\begin{grammarlens}[title={What you've built}]
One more word for asking for what I want.
\end{grammarlens}

\subsection*{Your turn}
\begin{itemize}
\raggedright
  \item \emph{Say it:} \emph{mǎidān}
  \item \emph{Say it:} \emph{mǎidān}, once more
  \item \emph{Say it:} \emph{fúwùyuán mǎidān} — Excuse me, the bill, please!
  \item \emph{Recall:} say the Mandarin for to roast, then the Mandarin for to fry, then say \emph{mǎidān} again
\end{itemize}

\begin{tcolorbox}[breakable,colback=teal!4,colframe=teal!35!black,title={Before you move on}]
What does mǎidān mean? (\textquotedblleft{}To pay the bill\textquotedblright{}.) Say it once more.
\end{tcolorbox}
\section[fù]{fù — to pay}

\label{lesson:ZH-C114-fu}



\begin{quote}
Before the new one: say the Mandarin for to fry, then the Mandarin for to pay the bill.
\end{quote}

\subsection*{You'll want to know: fù}
\emph{fù} — \textquotedblleft{}to pay\textquotedblright{}.

\begin{quote}
\emph{wǒ fù qián} — I'll pay.
\end{quote}

Its characters come later; for now it is a word for the ear.

\begin{cousinweb}
\emph{fù} is inherited from Middle Chinese \emph{pjuH} (Guangyun, in Baxter's transcription), glossed 'to give'.
\end{cousinweb}

\begin{grammarlens}[title={What you've built}]
One more word for asking for what I want.
\end{grammarlens}

\subsection*{Your turn}
\begin{itemize}
\raggedright
  \item \emph{Say it:} \emph{fù}
  \item \emph{Say it:} \emph{fù}, once more
  \item \emph{Say it:} \emph{wǒ fù qián} — I'll pay.
  \item \emph{Recall:} say the Mandarin for to fry, then the Mandarin for to pay the bill, then say \emph{fù} again
\end{itemize}

\begin{tcolorbox}[breakable,colback=teal!4,colframe=teal!35!black,title={Before you move on}]
What does fù mean? (\textquotedblleft{}To pay\textquotedblright{}.) Say it once more.
\end{tcolorbox}
\section[xiànjīn]{xiànjīn — cash}

\label{lesson:ZH-C114-xianjin}



\begin{quote}
Before the new one: say the Mandarin for to pay the bill, then the Mandarin for to pay.
\end{quote}

\subsection*{You'll want to know: xiànjīn}
\emph{xiànjīn} — \textquotedblleft{}cash\textquotedblright{}.

\begin{quote}
\emph{nǐ yǒu xiànjīn ma} — Have you got any cash?
\end{quote}

Its characters come later; for now it is a word for the ear.

\begin{cousinweb}
\emph{xiànjīn} is built from xiàn and jīn 'gold, money', both inherited from Middle Chinese, \emph{henH} and \emph{kim} (Guangyun, in Baxter's transcription). The dictionary glosses the first 'to show'; it is the xiàn of xiànzài 'now', and money there and then is the likely sense of the pair, though it is not settled.
\end{cousinweb}

\begin{grammarlens}[title={What you've built}]
One more word for asking for what I want.
\end{grammarlens}

\subsection*{Your turn}
\begin{itemize}
\raggedright
  \item \emph{Say it:} \emph{xiànjīn}
  \item \emph{Say it:} \emph{xiànjīn}, once more
  \item \emph{Say it:} \emph{nǐ yǒu xiànjīn ma} — Have you got any cash?
  \item \emph{Recall:} say the Mandarin for to pay the bill, then the Mandarin for to pay, then say \emph{xiànjīn} again
\end{itemize}

\begin{tcolorbox}[breakable,colback=teal!4,colframe=teal!35!black,title={Before you move on}]
What does xiànjīn mean? (\textquotedblleft{}Cash\textquotedblright{}.) Say it once more.
\end{tcolorbox}
\section[shuākǎ]{shuākǎ — to pay by card}

\label{lesson:ZH-C114-shuaka}



\begin{quote}
Before the new one: say the Mandarin for to pay, then the Mandarin for cash.
\end{quote}

\subsection*{You'll want to know: shuākǎ}
\emph{shuākǎ} — \textquotedblleft{}to pay by card\textquotedblright{}.

\begin{quote}
\emph{wǒ kěyǐ shuākǎ ma} — Can I pay by card?
\end{quote}

Its characters come later; for now it is a word for the ear.

kěyǐ is said kéyǐ: two third tones, and the first rises.

\begin{cousinweb}
\emph{shuākǎ} is built from shuā 'to brush, to swipe' and kǎ 'card'. The first is inherited from Middle Chinese \emph{srwaet} (Guangyun, in Baxter's transcription); the second is borrowed, usually traced to English \emph{card}.
\end{cousinweb}

\begin{grammarlens}[title={What you've built}]
One more word for asking for what I want.
\end{grammarlens}

\subsection*{Your turn}
\begin{itemize}
\raggedright
  \item \emph{Say it:} \emph{shuākǎ}
  \item \emph{Say it:} \emph{shuākǎ}, once more
  \item \emph{Say it:} \emph{wǒ kěyǐ shuākǎ ma} — Can I pay by card?
  \item \emph{Recall:} say the Mandarin for to pay, then the Mandarin for cash, then say \emph{shuākǎ} again
\end{itemize}

\begin{tcolorbox}[breakable,colback=teal!4,colframe=teal!35!black,title={Before you move on}]
What does shuākǎ mean? (\textquotedblleft{}To pay by card\textquotedblright{}.) Say it once more.
\end{tcolorbox}
\section[zhǎoqián]{zhǎoqián — to give change}

\label{lesson:ZH-C114-zhaoqian}



\begin{quote}
Before the new one: say the Mandarin for cash, then the Mandarin for to pay by card.
\end{quote}

\subsection*{You'll want to know: zhǎoqián}
\emph{zhǎoqián} — \textquotedblleft{}to give change\textquotedblright{}.

\begin{quote}
\emph{tā méiyǒu zhǎoqián} — He didn't give me my change.
\end{quote}

Its characters come later; for now it is a word for the ear.

\begin{cousinweb}
\emph{zhǎoqián} is built from zhǎo 'to look for' and qián 'money'. The second is inherited from Middle Chinese \emph{dzjen} (Guangyun, in Baxter's transcription); no dictionary we use records an older form of the first, so its history is not settled.
\end{cousinweb}

\begin{grammarlens}[title={What you've built}]
One more word for asking for what I want.
\end{grammarlens}

\subsection*{Your turn}
\begin{itemize}
\raggedright
  \item \emph{Say it:} \emph{zhǎoqián}
  \item \emph{Say it:} \emph{zhǎoqián}, once more
  \item \emph{Say it:} \emph{tā méiyǒu zhǎoqián} — He didn't give me my change.
  \item \emph{Recall:} say the Mandarin for cash, then the Mandarin for to pay by card, then say \emph{zhǎoqián} again
\end{itemize}

\begin{tcolorbox}[breakable,colback=teal!4,colframe=teal!35!black,title={Before you move on}]
What does zhǎoqián mean? (\textquotedblleft{}To give change\textquotedblright{}.) Say it once more.
\end{tcolorbox}
